
\subsubsection{Experimental Question}

\textbf{Does per-sample last-layer gradient cosine similarity with the within-batch mean gradient achieve higher ROC-AUC than per-sample loss as a predictor of spurious-minority group membership on Waterbirds and CelebA?}

Secondary questions: (1) How does the signal evolve across epochs? (2) Is failure consistent across datasets with different minority prevalences? (3) What does the gap magnitude suggest about the failure mechanism?

\subsubsection{Datasets}

\textbf{Waterbirds} [Wah et al., 2011; Sagawa et al., 2020]: Synthetic spurious correlation (bird species $\times$ background, 95\% strength). Train: 4,795 samples, ~82\% majority. Minority groups for evaluation: landbird on water and waterbird on land (group IDs {1, 2} in the group\_DRO encoding).

\textbf{CelebA} \citep{Liu2015}: Blond hair prediction; spurious attribute: gender. Train: 16,000 samples (stratified subsample, ~10\% of full 162K dataset, seed 42). Minority: blond male (~0.8\% of training data, group ID 3).

CelebA complements Waterbirds with more extreme minority prevalence --- key for testing whether prevalence modulates the contamination effect.

\subsubsection{Baselines}

\begin{itemize}
\item \textbf{Gradient Alignment Score ($a_i$):} Negated within-batch cosine similarity (Section 3)
\item \textbf{Per-Sample Loss ($\ell_i$):} Cross-entropy loss; signal used by JTT and LfF
\end{itemize}

Identical model states; the single variable is signal type. We evaluate signal discriminability (ROC-AUC) rather than downstream worst-group accuracy; the latter is only relevant if the signal exists.

\subsubsection{Evaluation}

\textbf{ROC-AUC} (sklearn): threshold-free discriminability measure. AUC < 0.5 indicates inverted signal. Binary labels: $y_i = \mathbb{1}[\text{group\_id}_i \in \text{minority}]$.

